\section{Objectifs}
\begin{itemize}
\item Calculer une probabilité conditionnelle.
\item Distinguer $P(A\cap B)$, $P_A(B)$ et $P_B(A)$.
\item Comprendre l'indépendance de deux évènements.
\item Représenter de courtes répétitions de Bernoulli.
\end{itemize}

\section{1. Probabilité conditionnelle}
Lorsque $P(A)>0$, la probabilité de $B$ sachant $A$ est
\[
P_A(B)=\frac{P(A\cap B)}{P(A)}.
\]
Ainsi,
\[
P(A\cap B)=P(A)\times P_A(B).
\]

Sur un arbre pondéré, la probabilité d'un chemin s'obtient en multipliant les probabilités de ses branches.

\section{2. Bien distinguer les notations}
\begin{itemize}
\item $P(A\cap B)$ : A et B se réalisent.
\item $P_A(B)$ : probabilité de B lorsque A est réalisé.
\item $P_B(A)$ : probabilité de A lorsque B est réalisé.
\end{itemize}

Ces trois nombres n'ont aucune raison d'être égaux.

\section{3. Indépendance}
Deux évènements A et B sont indépendants si la connaissance de l'un ne modifie pas la probabilité de l'autre. On peut utiliser
\[
P(A\cap B)=P(A)P(B).
\]

Indépendance ne signifie pas incompatibilité. Deux évènements incompatibles de probabilités non nulles ne peuvent d'ailleurs pas être indépendants.

\section{4. Répétitions de Bernoulli}
Une épreuve de Bernoulli possède deux issues : succès et échec. Le programme demande de savoir représenter la répétition de $n$ épreuves identiques et indépendantes avec $n\le4$.

Pour trois répétitions de probabilité de succès $p=0{,}7$, la probabilité d'obtenir exactement deux succès est
\[
3\times0{,}7^2\times0{,}3=0{,}441.
\]

La formule générale de la loi binomiale n'est pas nécessaire dans ce module.

\section{5. Modèle et réalité}
Le modèle probabiliste est une hypothèse. Une simulation permet de comparer des fréquences observées aux probabilités du modèle, mais ne transforme pas une hypothèse en certitude.

\section{Automatismes à entretenir}
\begin{itemize}
\item Lire un tableau croisé.
\item Calculer une fréquence conditionnelle.
\item Multiplier les probabilités le long d'un chemin.
\item Distinguer intersection et conditionnement.
\end{itemize}

\section{Exemple : test médical}
Supposons que 2 \% d'une population soit malade et qu'un test soit positif chez 90 \% des malades :
\[
P(M\cap +)=0{,}02\times0{,}90=0{,}018.
\]
Cette valeur n'est pas la probabilité d'être malade sachant que le test est positif.

\section{Erreurs fréquentes}
\begin{itemize}
\item Confondre $P_A(B)$ et $P_B(A)$.
\item Dire que deux évènements indépendants sont incompatibles.
\item Additionner au lieu de multiplier le long d'un chemin.
\item Introduire une loi binomiale générale qui n'est pas attendue ici.
\end{itemize}

\section{À retenir}
Le conditionnement décrit une probabilité sous information. L'indépendance signifie que cette information ne change pas la probabilité de l'évènement considéré.
